\documentclass[11pt]{article}
\usepackage[margin=1in]{geometry}
\usepackage{graphicx,booktabs,float,xcolor,hyperref}
\usepackage{helvet}\renewcommand{\familydefault}{\sfdefault}
\hypersetup{colorlinks=true,linkcolor=blue,urlcolor=blue}
\newcommand{\Mw}{$M_{\mathrm w}$}
\setlength{\parindent}{0pt}\setlength{\parskip}{6pt}
\title{\textbf{Prospective source model of the 2026 \Mw~7.3 Mexico (Chiapas) earthquake}\\
\large A blind test of our MTUQ$+$WISP pipeline, produced before the USGS finite fault}
\author{Automated pipeline (17.Venezuela\_2026)}
\date{Frozen 2026-07-18, hours after origin --- for later comparison with the USGS solution}
\begin{document}\maketitle

\section*{Purpose}
The USGS had not yet posted a finite-fault model for \texttt{us7000t1bu} when this was produced.
This is a \textbf{prospective (blind) test}: we run our teleseismic MTUQ moment-tensor $+$ WISP
finite-fault pipeline on open GSN data and freeze the prediction here so it can be compared with the
USGS solution once posted. Everything below uses \emph{only} teleseismic body $+$ surface waves and
the USGS W-phase moment tensor as the seed; \textbf{no InSAR, regional, or geodetic data} (which
USGS will likely add, and which --- per our Al Haouz study --- tends to yield a smaller, more
compact source).

\section{Event and USGS preliminary benchmark (frozen)}
\begin{table}[H]\centering\small
\begin{tabular}{ll}
\toprule
Origin & 2026-07-17T14:48:39.77 UTC, 14.604$^\circ$N, 92.952$^\circ$W, depth 18.6 km \\
Magnitude & \Mw~7.3 (W-phase); $M_0 = 1.1\times10^{20}$ N\,m; centroid depth 15.5 km \\
Tectonics & Middle America subduction zone (Cocos plate), \textbf{shallow thrust} \\
W-phase NP1 (steep) & strike 122.75, dip 63.86, rake 97.74 \\
W-phase NP2 (shallow) & strike 285.6, dip 27.19, rake 74.65 \quad(likely megathrust interface) \\
\bottomrule
\end{tabular}
\end{table}

This is analogous in size to the 2024/2026 Atacama and 2026 Venezuela events we have modelled, but
it is a \emph{shallow subduction thrust}, so we apply the shallow-source lessons from our Al Haouz
(Morocco) study.

\section{Data}
23 GSN broadband stations (networks IU, II, G, GE, \ldots) at 31--86$^\circ$, fetched from IRIS.
WISP uses 15 with usable P; MTUQ uses 11 three-component stations. The azimuthal coverage has a
$\sim$87$^\circ$ gap over the Pacific (few island stations) --- adequate but not ideal.

\begin{figure}[H]\centering
\includegraphics[width=0.7\textwidth]{../fig_mexico_stations.png}
\caption{Teleseismic station distribution (azimuthal-equidistant, rings at 30/60/90$^\circ$).}
\end{figure}

\section{Moment tensor (MTUQ) --- the shallow-thrust trap, a third time}
For this shallow source the same pattern we documented on Al Haouz reappears, now at \Mw~7.3:

\begin{table}[H]\centering\small
\begin{tabular}{llll}
\toprule
MTUQ solution & strike/dip/rake & \Mw & verdict \\
\midrule
\textbf{DC, body-only (adopted)} & \textbf{284/30/79} & \textbf{7.30} & $\approx$ W-phase NP2 (shallow interface) \\
DC, surface waves & 121/85/88 & 7.45 & dip 85$^\circ$ (too steep --- shallow indeterminacy) \\
full moment tensor & 126/81/81 & 7.60 & destabilised ($w=0.90$, huge non-DC) \\
W-phase (USGS) & 122.75/63.86/98 \;/\; 285.6/27.19/75 & 7.3 & reference \\
\bottomrule
\end{tabular}
\end{table}

The \textbf{body-only DC recovers the W-phase mechanism and moment} (284/30 matches the shallow NP2;
\Mw~7.30 matches $M_{ww}$). The surface-wave DC pushes the dip to 85$^\circ$ (the Kanamori--Given
shallow dip-slip indeterminacy), and the full moment tensor is destabilised into a large spurious
non-double-couple with \Mw~7.60. This is the third independent confirmation of our rule:
\textbf{for a shallow source, invert body waves for the mechanism.} We adopt the body-only DC.

\begin{figure}[H]\centering
\includegraphics[width=0.42\textwidth]{../../results/mtuq_dc_2026mexico/2026mexico_beachball.png}
\caption{Adopted body-only DC moment tensor (284/30/79) with the station distribution --- a shallow
subduction thrust, consistent with the W-phase.}
\end{figure}

\section{Finite fault (WISP), both nodal planes}
WISP \texttt{auto\_model} on both W-phase nodal planes, \textbf{seeded from the W-phase moment}
(the key lesson from our Al Haouz ablation: seed from the authoritative $M_{ww}$). Teleseismic body
$+$ surface waves. We then applied the \textbf{cross-correlation time-shift step}
(\texttt{shift\_match}; \emph{not} run by \texttt{auto\_model} --- a separate manual step) and
re-inverted, which cut the misfit \textbf{13--16\%} on both planes with \Mw\ unchanged (as on Al
Haouz). Values below are the shifted models.

\begin{table}[H]\centering\small
\begin{tabular}{llllll}
\toprule
plane & strike/dip & \Mw & misfit (pre$\to$post shift) & $L\times W$ / peak slip & slip depth / T$_{5-95}$ \\
\midrule
NP1 steep & 123/64 & 7.29 & 0.1120 $\to$ \textbf{0.0940} & 82$\times$41 km / 3.43 m & 1--34 km / 19.2 s \\
\textbf{NP2 shallow (interface)} & \textbf{286/27} & \textbf{7.29} & 0.1114 $\to$ \textbf{0.0969} & 82$\times$68 km / 2.54 m & 3--24 km / 20.3 s \\
\bottomrule
\end{tabular}
\caption{Both planes fit almost identically --- teleseismic waveforms are a weak plane discriminant,
as we found for Al Haouz and Chile. For a subduction thrust the shallow-dipping NP2 is the megathrust
interface and fits marginally better; we \emph{prefer NP2} but show both. Both give \Mw~7.29
(matching $M_{ww}$), a compact updip asperity, and $\sim$20 s duration. The time-shift step improved
each fit $\sim$13--16\% and slightly sharpened the slip (NP2 peak 2.29$\to$2.54 m).}
\end{table}

\begin{figure}[H]\centering
\includegraphics[width=0.88\textwidth]{../20260717144839/ffm.0/NP2_shift/plots/SlipIsochrone_trueScale.png}
\caption{Preferred (NP2, shallow interface) on-fault slip at true 1:1 scale; star $=$ hypocentre,
white contours $=$ rupture front every 5 s, arrows $=$ thrust slip. Compact updip asperity, peak
2.5 m, slip $\sim$5--24 km depth (shifted model).}
\end{figure}
\begin{figure}[H]\centering
\includegraphics[width=0.55\textwidth]{../20260717144839/ffm.0/NP2_shift/plots/MomentRate.png}
\caption{Moment-rate function (NP2): peak at $\sim$11 s, total duration $\sim$20 s.}
\end{figure}
\begin{figure}[H]\centering
\includegraphics[width=0.92\textwidth]{../20260717144839/ffm.0/NP2_shift/plots/Map.png}
\caption{Map view of the NP2 slip projected to the surface.}
\end{figure}
\begin{figure}[H]\centering
\includegraphics[width=0.92\textwidth]{../20260717144839/ffm.0/NP2_shift/plots/P_body_waves.png}
\caption{WISP teleseismic P-wave fits (NP2). Black observed, red synthetic.}
\end{figure}

\section{Prospective prediction (compare against the USGS FFM when posted)}
\begin{itemize}
\item \textbf{Mechanism:} shallow subduction thrust; preferred plane $\approx$285/27/75 (megathrust
interface), or steep conjugate 123/64.
\item \textbf{\Mw:} 7.29--7.30 ($M_0 \approx 1.1\times10^{20}$ N\,m), matching $M_{ww}$.
\item \textbf{Peak slip:} $\sim$2.5--3.4 m; \textbf{slip depth} $\sim$5--25 km; single compact updip
asperity near/above the hypocentre.
\item \textbf{Duration:} $\sim$20 s, moment-rate peak at $\sim$10 s.
\end{itemize}

\textbf{Honest caveats.} (1) Teleseismic-only: no InSAR/regional yet. Per our Al Haouz study, adding
geodesy usually shrinks the source and raises peak slip. (2) Location is the catalogue hypocentre
(not relocated to a geodetic centroid). (3) The $\sim$87$^\circ$ Pacific azimuthal gap limits
resolution on the down-dip/along-strike updip direction.

\section{Comparison with the USGS finite fault (posted after our freeze)}
USGS (D.~Goldberg) posted their finite fault on 2026-07-17. Crucially it is \textbf{also
teleseismic-only} (BODY 63, SURF 52, InSAR 0) --- a true like-for-like comparison.

\begin{table}[H]\centering\small
\begin{tabular}{lll l}
\toprule
 & USGS (Goldberg) & Ours (NP2\_shift) & assessment \\
\midrule
Fault plane & 286/27/75 & 286/27/75 & \textbf{exact match} \\
\Mw & 7.29 & 7.29 & exact \\
Peak slip & 2.85 m & 2.54 m & close ($-$11\%) \\
Slip depth ($>$15\%) & 4--25 km & 3--24 km & \textbf{excellent match} \\
Moment-rate peak & 10.2 s & 10.2 s & exact \\
Data used & teleseismic-only & teleseismic-only & match \\
Fault template & 125$\times$65 km & 82$\times$68 km & USGS larger \\
Rupture velocity & 1.85 km/s & 2.08 km/s & ours faster \\
Duration T$_{5-95}$ & 35 s & 20 s & USGS longer \\
\bottomrule
\end{tabular}
\caption{Blind-test outcome. The source \emph{fundamentals} agree closely: the same nodal plane
(we correctly preferred the shallow megathrust interface), the same \Mw, the same slip depth, the
same peak-slip magnitude, and the same main-pulse timing --- both teleseismic-only. The differences
are secondary modelling choices: USGS used a larger fault template (125 km) with a slower rupture
velocity, producing a longer, more spread-out rupture; ours is more compact.}
\end{table}

\begin{figure}[H]\centering
\includegraphics[width=\textwidth]{../fig_usgs_comparison.png}
\caption{Left: moment-rate functions --- the main pulse (which carries most of the moment) is nearly
identical in timing and amplitude; USGS spreads the last $\sim$15\% of moment over a longer tail
(bigger fault $+$ slower V$_r$). Right: slip-vs-depth --- the distributions nearly overlap, both
peaking at $\sim$16 km.}
\end{figure}

\textbf{Where our forecast missed (honest).} Our two caveats predicted the \emph{opposite} of what
USGS did: we expected them to add InSAR and produce a \emph{more compact} source. Instead their
model is teleseismic-only and \emph{longer} than ours. So our prediction was, if anything,
conservatively compact --- the mechanism, moment, slip depth and main-pulse timing all validated,
while the total rupture length/duration is a genuine model-choice difference (larger template $+$
slower V$_r$) on which the teleseismic data are weakly constraining.

\end{document}
